\section{Endpoint analysis for wrapped EC}
\label{app:analytic-details}

This appendix proves the sparse and dense claims used in
Lemma~\ref{lem:ec-endpoints}.  It also explains why one set of constants
controls a whole interval of input probabilities.  Fix $m\geq1$ and write
$\varrho_m(a,z)$ for the largest absolute eigenvalue of
$\mathbf N_{m,a}(z)$.

At $(a,z)=(0,1)$, the reduced state chain has two closed classes.  The active
class consists of states $0,\ldots,m-1$, and state $m$ is the all-zero class.
Define
\begin{equation}
 c_m:=\frac{1}{2^m-1},\qquad
 \alpha_m:=\frac{2^{m-1}}{2^m-1}.
 \label{eq:ec-class-parameters}
\end{equation}
Within the active class, the stationary probability of state $j$ is
$\alpha_m2^{-j}$.  Thus $c_m$ is the probability of the boundary state
$m-1$, and $\alpha_m$ is the mean emitted weight in one active step.

Set $z=e^{-\theta a}$ for $\theta\geq0$.  Restricting the first-order
perturbation of $\mathbf N_{m,a}(e^{-\theta a})$ to the two closed classes
gives
\begin{equation}
 \mathbf G_m(\theta):=
 \begin{pmatrix}
  -c_m-\theta\alpha_m&c_m\\
  1&-1
 \end{pmatrix}.
 \label{eq:ec-effective-generator}
\end{equation}
The first row represents the active class and the second the all-zero class.
The off-diagonal entries are the first-order rates at which a small input
probability moves between the classes.  The marker contributes the additional
rate $\theta\alpha_m$ while the chain is active.

The larger eigenvalue of \eqref{eq:ec-effective-generator} is
\begin{equation}
 g_m(\theta)=
 -\frac{1+c_m+\theta\alpha_m}{2}
 +\frac{\sqrt{(1-c_m-\theta\alpha_m)^2+4c_m}}{2}.
 \label{eq:ec-effective-eigenvalue}
\end{equation}

The next lemma supplies the uniform remainder needed by the sparse-message
argument.  It is stated for a bounded marker parameter because the proof of
Lemma~\ref{lem:ec-endpoints} fixes one value of $\theta$.

\begin{lemma}[Two-class perturbation]
\label{lem:ec-two-class-perturbation}
Fix $m\geq1$ and $\Theta<\infty$.  There are constants
$a_0,K>0$, depending only on $m$ and $\Theta$, such that
\begin{equation}
 \left|
 \ln\varrho_m(a,e^{-\theta a})-a g_m(\theta)
 \right|\leq Ka^2
 \label{eq:ec-uniform-perturbation}
\end{equation}
for $0<a\leq a_0$ and $0\leq\theta\leq\Theta$.
The positive Perron eigenvector can be normalized so that the ratio of its
largest coordinate to its smallest coordinate is at most $K$ on the same
domain.
\end{lemma}

\begin{proof}
Put $\mathbf P(a,\theta):=\mathbf N_{m,a}(e^{-\theta a})$.
Its entries are analytic in $(a,\theta)$ near
$\{0\}\times[0,\Theta]$.  After decreasing $a_0$ if necessary, Taylor's
theorem on the resulting compact set gives
\begin{equation}
 \mathbf P(a,\theta)=\mathbf P_0+a\mathbf P_1(\theta)
 +\mathbf E(a,\theta),
 \qquad \|\mathbf E(a,\theta)\|\leq K_0a^2,
 \label{eq:ec-matrix-taylor}
\end{equation}
for a constant $K_0$ and any fixed matrix norm.

The eigenvalue one of $\mathbf P_0$ is semisimple with multiplicity two.
Its two right eigenspaces are the active and all-zero closed classes.
All remaining eigenvalues have modulus less than one.  Let $\mathbf U$
contain one right eigenvector for each closed class, and let $\mathbf V$
contain the dual stationary left eigenvectors, normalized by
$\mathbf V^{\mathsf T}\mathbf U=\mathbf I_2$.  Direct substitution of the
transitions in \eqref{eq:ec-ber-active}--\eqref{eq:ec-ber-zero} gives
\begin{equation}
 \mathbf V^{\mathsf T}\mathbf P_1(\theta)\mathbf U
 =\mathbf G_m(\theta).
 \label{eq:ec-compressed-derivative}
\end{equation}

Use the spectral projection of $\mathbf P_0$ to decompose the state space
into the two-dimensional one-eigenspace and its invariant complement.  In
this basis, the lower-right block of
$\mathbf P_0-\lambda\mathbf I$ is uniformly invertible whenever
$|\lambda-1|$ is sufficiently small.  The Schur complement of that block,
together with \eqref{eq:ec-matrix-taylor}, is
\begin{equation}
 (1-\lambda)\mathbf I_2
 +a\mathbf G_m(\theta)+\mathbf R(a,\theta,\lambda),
 \qquad \|\mathbf R(a,\theta,\lambda)\|\leq K_1a^2.
 \label{eq:ec-schur-complement}
\end{equation}
The constants are uniform for $\theta\in[0,\Theta]$.  Therefore the two
eigenvalues near one are
$1+a\mu_i(\theta)+O(a^2)$, where the $\mu_i(\theta)$ are the eigenvalues of
$\mathbf G_m(\theta)$.  The remaining eigenvalues stay a fixed distance
inside the unit circle.

For $a>0$ and $e^{-\theta a}>0$, the nonnegative matrix
$\mathbf P(a,\theta)$ is irreducible.  Its Perron eigenvalue is therefore the
branch associated with the larger eigenvalue $g_m(\theta)$.  The two
eigenvalues of $\mathbf G_m(\theta)$ are distinct throughout the compact
interval because $c_m>0$.  The Schur-complement equation also shows that the
normalized Perron eigenvector converges uniformly to the lift of the positive
eigenvector of $\mathbf G_m(\theta)$.  Its coordinates are consequently
bounded above and away from zero.  Finally,
$\ln(1+a g_m(\theta)+O(a^2))=a g_m(\theta)+O(a^2)$, with a uniform constant.
This proves both claims.
\end{proof}

Lemma~\ref{lem:ec-two-class-perturbation} gives
\begin{equation}
 \ln\varrho_m(a,e^{-\theta a})
 =a g_m(\theta)+O(a^2).
 \label{eq:ec-sparse-expansion}
\end{equation}

Substituting this marker into \eqref{eq:ec-rate-function} gives
\begin{equation}
 \liminf_{a\downarrow0}\frac{J^{\mathrm w}_{m,\delta}(a)}a
 \geq\sup_{\theta\geq0}\{-\delta\theta-g_m(\theta)\}.
 \label{eq:ec-effective-optimization}
\end{equation}
Put $x:=\delta/\alpha_m$.  Direct one-variable optimization yields
\begin{align}
 \sup_{\theta\geq0}\{-\delta\theta-g_m(\theta)\}
 &=\left(\sqrt{1-x}-\sqrt{c_mx}\right)^2\\
 &=\left(
 \sqrt{1-2\delta+2^{1-m}\delta}
 -\sqrt{2^{1-m}\delta}
 \right)^2
 =I^{\mathrm w}_{m,\delta}.
 \label{eq:ec-effective-rate}
\end{align}
This proves the sparse exponent in \eqref{eq:ec-sparse-endpoint}.

We next justify the uniform probability bound that the asymptotic distance
proof needs.  For a positive matrix and a positive right eigenvector $v$ with
eigenvalue $\lambda$, componentwise comparison gives
\begin{equation}
 \e_m^{\mathsf T}\mathbf N_{m,a}(z)^n\one
 \leq \frac{\max_i v_i}{\min_i v_i}\lambda^n.
 \label{eq:ec-perron-prefactor}
\end{equation}
For a fixed $\theta$ below the sparse threshold,
Lemma~\ref{lem:ec-two-class-perturbation} bounds the coordinate ratio of the
Perron eigenvector for all sufficiently small $a$.  Equations
\eqref{eq:ec-sparse-expansion} and \eqref{eq:ec-perron-prefactor}, together
with the marker inequality, give constants $K,a_0>0$ such that
\begin{equation}
 \Pr[H_{n,a}\leq\lfloor\delta n\rfloor]
 \leq K\exp\{-na(I^{\mathrm w}_{m,\delta}-\varepsilon)\}
 \label{eq:ec-uniform-sparse-tail}
\end{equation}
for all $n\geq1$ and $a\in(0,a_0]$.

For input probabilities bounded away from zero, the matrix at $z=1$ is a
finite irreducible Markov transition matrix.  Its stationary law satisfies
\begin{equation}
 \pi_j=\pi_0 2^{-j}\ (1\leq j<m),\qquad
 \pi_m=\pi_{m-1},\qquad \pi_0=1/2.
 \label{eq:ec-stationary-law}
\end{equation}
Consequently, the derivative of $\ln\varrho_m(a,z)$ with respect to $\ln z$
at $z=1$ is $1/2$.  Because $\delta<1/2$, a marker slightly below one gives
a positive lower-tail exponent.  Continuity and a finite cover of every
closed interval $[a_*,1/2]$ give constants that are uniform throughout that
interval.

Finally, when $a=1/2$, every row of $\mathbf N_{m,1/2}(z)$ sums to
$(1+z)/2$.  Hence
\begin{equation}
 \varrho_m(1/2,z)=\frac{1+z}{2}.
\end{equation}
Optimizing $\delta\ln z-\ln((1+z)/2)$ gives
$\ln2-h(\delta)$ and proves \eqref{eq:ec-dense-endpoint}.
